\documentclass[conference]{IEEEtran}
\IEEEoverridecommandlockouts
\usepackage{cite}
\usepackage{amsmath,amssymb,amsfonts}
\usepackage{algorithmic}
\usepackage{graphicx}
\usepackage{textcomp}
\usepackage{xcolor}
\usepackage{booktabs}
\usepackage{multirow}
\usepackage{hyperref}

\def\BibTeX{{\rm B\kern-.05em{\sc i\kern-.025em b}\kern-.08em
    T\kern-.1667em\lower.7ex\hbox{E}\kern-.125emX}}

\begin{document}

\title{Training-Free Multimodal Image Fusion by Anisotropic-Diffusion Decomposition and Local-Statistics Weighting: An Ablation-Driven Study}

\author{\IEEEauthorblockN{Advanced Visual Computing \& Multimodal Imaging Laboratory}
\IEEEauthorblockA{\textit{Department of Computer Science and Engineering} \\
\textit{Autonomous Vision and Image Processing Group}\\
Email: research-draft@lab.org}
}

\maketitle

\begin{abstract}
Multimodal image fusion combines complementary information from distinct sensors into a single unified composite image. Recent deep neural architectures report strong results but require extensive training datasets and compute, and several are evaluated on benchmark datasets that overlap directly with their own training corpora. We investigate a deterministic, training-free alternative: the luminance channel is decomposed by Perona-Malik anisotropic diffusion into a gross (base) layer and a fine (detail) layer, and the two layers are fused via saliency weights computed from local statistics and refined by guided filtering. Color is preserved by fusing luminance only in YCbCr space. We evaluate 14 design choices across 120 development pairs using an objective selection rule fixed prior to the final sweep, and evaluate the frozen variant on 683 held-out test pairs across 7 public datasets (MSRS, RoadScene, TNO, CT-MRI, PET-MRI, SPECT-MRI, Lytro) against 5 classical and 4 deep baselines. The development study shows that eliminating edge compensation ($\gamma=0.0$) and operating without feature clustering ($K=1$) yields the top-ranked configuration (Mean Rank 2.8 vs 13.6 for prototype), reducing artifact noise ($N^{AB/F}$) from 0.2888 to 0.0153. On held-out test data, the method achieves competitive structural fidelity ($SSIM = 0.7675$, $Q^{AB/F} = 0.6296$ on CT-MRI) with a per-pair CPU latency of 88.2 ms on an Intel Core i5-12450H CPU. While GFF achieves slightly higher edge transfer on high-frequency boundaries, our pipeline produces significantly lower artifact noise and superior mutual information without learned parameters.
\end{abstract}

\begin{IEEEkeywords}
image fusion, anisotropic diffusion, guided filtering, training-free, multimodal imaging, ablation.
\end{IEEEkeywords}

\section{Introduction}
\subsection{Motivation}
Multimodal sensor fusion plays an indispensable role in medical diagnosis, autonomous navigation, and optical surveillance. In clinical oncology and neurology, Magnetic Resonance Imaging (MRI) provides exceptional soft-tissue differentiation, whereas Computed Tomography (CT) captures dense bone morphology \cite{ref_harvard}. Positron Emission Tomography (PET) and Single-Photon Emission Computed Tomography (SPECT) provide physiological metabolic activity \cite{ref_harvard}. In transportation and night vision, visible cameras capture rich chromatic texture while infrared (IR) cameras register thermal radiation \cite{ref_tno, ref_piafusion}. Multi-focus optical imaging integrates shallow depth-of-field captures into an all-in-focus composite \cite{ref_lytro}.

\subsection{Problem Formulation}
Deep learning fusion models demand large-scale training data, lack verifiable physical constraints, and frequently suffer from test-set contamination where published metrics are evaluated on the exact datasets used for network training (audited in Section IV-C). Classical methods are training-free but often distort sharp boundaries or inject high-frequency noise. We investigate how far a simple, interpretable diffusion-based pipeline can go, and systematically dissect which architectural components genuinely contribute.

\subsection{Contributions}
\begin{enumerate}
    \item \textbf{Deterministic Pipeline}: A training-free fusion architecture combining Perona-Malik anisotropic diffusion, local-statistics weighting, and guided filtering with zero learned parameters.
    \item \textbf{Pre-Registered Ablation}: A 120-pair development ablation showing that edge compensation ($\gamma > 0$) and clustering ($K > 1$) actively harm structural fidelity.
    \item \textbf{Comprehensive Test Benchmark}: Evaluation across 683 pairs from 7 datasets against 5 classical and 4 deep baselines with full disclosure of baseline training data overlap.
    \item \textbf{Reproducibility}: Complete release of code, parameter configurations, split manifests, and per-image evaluation tables.
\end{enumerate}

\section{Related Work}
\subsection{Multiscale and Edge-Preserving Methods}
Discrete Wavelet Transforms (DWT) and Laplacian pyramids frequently suffer from ringing artifacts around high-contrast edges. Edge-preserving filters, such as Guided Filter Fusion (GFF, Li et al. \cite{ref_gff}), Cross Bilateral Filters (CBF, Kumar \cite{ref_cbf}), and Anisotropic Diffusion Fusion (ADF, Bavirisetti \& Dhuli \cite{ref_bavirisetti}), alleviate these artifacts. Unlike Gaussian or wavelet pyramids, our pipeline uses nonlinear Perona-Malik diffusion to form a strictly additive two-layer decomposition ($GT + FT = Y$) without downsampling and refines saliency weights with guided filtering using the source luminance as structural guide.

\subsection{Deep Methods \& Training Data Audit}
Deep architectures include DenseFuse \cite{ref_densefuse} (convolutional autoencoder), U2Fusion \cite{ref_u2fusion} (DenseNet with continuous similarity loss), SwinFusion \cite{ref_swinfusion} (cross-domain transformer), and CDDFuse \cite{ref_cddfuse} (Restormer with dual-branch decomposition). Table II provides a critical disclosure: while DenseFuse was trained on MS-COCO natural images, U2Fusion, SwinFusion, and CDDFuse were trained directly on the benchmark datasets (MSRS, RoadScene, Harvard Whole Brain Atlas) on which they are evaluated.

\section{Methodology}
\subsection{Color Handling}
Inputs are converted to ITU-R BT.601 YCbCr space:
\begin{equation}
Y = 0.299 R + 0.587 G + 0.114 B
\end{equation}
Fusion operates strictly on luminance $Y \in [0, 1]$. Chrominance channels ($Cb, Cr$) are preserved from the color source (or averaged for dual-color inputs) and inverted back to RGB, eliminating color cast distortion.

\subsection{Anisotropic Diffusion \& Numerical Stability}
Perona-Malik diffusion decomposes luminance $u$ via:
\begin{equation}
\frac{\partial u}{\partial t} = \nabla \cdot ( g(\|\nabla u\|) \nabla u )
\end{equation}
with exponential conduction $g(s) = \exp(-(s/\kappa)^2)$, where $\kappa = 0.1176$. On a 4-neighborhood discrete grid:
\begin{equation}
u^{t+1} = u^t + \lambda \sum_{d \in \{N,S,E,W\}} g(|\nabla_d u^t|) \nabla_d u^t
\end{equation}
Because $g \in (0, 1]$, choosing $\lambda = 0.20$ strictly satisfies the CFL stability condition ($\lambda \le 0.25$), guaranteeing numerical boundedness and maximum principle preservation.

\subsection{Two-Scale Decomposition \& Saliency Weighting}
After $T=20$ iterations, the gross (base) layer is $GT = u(T)$ and the fine (detail) layer is $FT = Y - GT$, ensuring $GT + FT = Y$ holds exactly by construction. Local saliency $S_i(x)$ is estimated within a $7 \times 7$ window using local variance and spatial frequency:
\begin{equation}
S_i(x) = \text{BoxFilter}((Y_i - \mu_i)^2, w=7) + \sqrt{RF_i^2 + CF_i^2}
\end{equation}
Initial weights $W_i = S_i / (\sum S_j + \epsilon)$ are regularized via guided filtering with luminance $Y$ as the guide (radius $r=5$, $\epsilon=10^{-4}$), yielding edge-aligned fusion maps $F_{\text{layer}} = \sum_i W_i L_i$.

\subsection{Reconstruction \& The Edge-Compensation Fallacy}
In the frozen variant CADM-MSWN v2, composite luminance is reconstructed cleanly as:
\begin{equation}
F = F_{GT} + F_{FT}
\end{equation}
The preliminary prototype added an edge-compensation term $\gamma \cdot (FE_1 + FE_2)$ with $\gamma=0.5$. Because fine-edge maps represent gradient magnitudes rather than intensity flux, directly adding them to pixel intensities mixed quantities of disparate physical units, injecting severe gradient ringing and elevating $N^{AB/F}$ to 0.2888. Setting $\gamma=0$ completely eliminates this artifact.

\subsection{Feature Clustering \& Complexity}
Our development sweep tested k-means clustering with $K \in \{1, 2, 3\}$. Unclustered global weighting ($K=1$) achieved Rank 1 (Mean Rank 2.8), whereas $K=2$ (Rank 3) and $K=3$ (Rank 6) fragmented continuous boundary weights. Total computational complexity is strictly $\mathcal{O}(T \cdot H \cdot W)$ for diffusion, $\mathcal{O}(H \cdot W)$ for box-filtered local statistics, and $\mathcal{O}(H \cdot W)$ for guided filtering, yielding deterministic $\mathcal{O}(H \cdot W)$ overall scaling.

\section{Experimental Protocol}
\subsection{Datasets and Split}
We evaluate on 120 development pairs and 683 test pairs across 7 public datasets. Excluded early inspection pairs: CT-MRI 16003, PET-MRI 25015, SPECT-MRI 10007, lytro\_01, and RoadScene FLIR\_00006.

\begin{table}[htbp]
\caption{Multimodal Benchmark Dataset Inventory and Split Structure.}
\begin{center}
\begin{tabular}{llccl}
\toprule
\textbf{Dataset} & \textbf{Modality} & \textbf{Dev} & \textbf{Test} & \textbf{Resolution} \\
\midrule
MSRS & IR-VIS & 100 & 361 & 480$\times$640 \\
RoadScene & IR-VIS & 0 & 220 & Various \\
TNO & IR-VIS & 0 & 42 & Various \\
Harvard CT-MRI & Medical & 5 & 15 & 256$\times$256 \\
Harvard PET-MRI & Medical & 5 & 15 & 256$\times$256 \\
Harvard SPECT-MRI & Medical & 5 & 15 & 256$\times$256 \\
Lytro & Multi-focus & 5 & 15 & 520$\times$520 \\
\midrule
\textbf{Total} & & \textbf{120} & \textbf{683} & \\
\bottomrule
\end{tabular}
\end{center}
\end{table}

\subsection{Selection Protocol}
Design choices were compared on the dev set with a one-at-a-time sweep and ranked by mean rank over 5 pre-specified primary metrics ($Q^{AB/F}, N^{AB/F}, MI, SSIM, VIFP$). Ties broken by latency. Hyperparameters frozen in \texttt{config.yaml} before test execution.

\subsection{Compared Methods \& Training Data Disclosure}
\begin{table}[htbp]
\caption{Deep Baseline Training Data Audit and Domain Overlap.}
\begin{center}
\begin{tabular}{lll}
\toprule
\textbf{Method} & \textbf{Training Data} & \textbf{Overlap with Test Sets} \\
\midrule
DenseFuse \cite{ref_densefuse} & MS-COCO (80k images) & Zero (natural domain) \\
U2Fusion \cite{ref_u2fusion} & RoadScene, KAIST, Harvard & High Overlap (in-domain) \\
SwinFusion \cite{ref_swinfusion} & MSRS, Harvard, Real MFF & High Overlap (in-domain) \\
CDDFuse \cite{ref_cddfuse} & MSRS, Harvard Atlas & High Overlap (in-domain) \\
\bottomrule
\end{tabular}
\end{center}
\end{table}

\subsection{Metrics}
Primary metrics: $Q^{AB/F}$ \cite{ref_qabf}, $N^{AB/F}$ \cite{ref_nabf}, $MI$, $SSIM$ \cite{ref_ssim}, and $VIFP$ \cite{ref_vif}. Secondary metrics: $EN$, $SD$, $SF$, $AG$. Metric implementations validated against external toolboxes with documented mathematical consistency.

\section{Results}
\subsection{Development-Set Ablation Analysis}
Table III documents the complete 14-configuration evaluation on all 120 dev pairs across the unified metric suite. $K=1$ without edge compensation ($\gamma=0.0$) placed Rank 1 overall with Mean Rank 2.8, while the prototype ($\gamma=0.5$) placed Rank 14 (Mean Rank 13.6).

\begin{table*}[htbp]
\caption{Pre-Registered Development-Set Ablation Sweep across 120 Pairs (14 Configurations).}
\begin{center}
\begin{tabular}{clccccccccccc}
\toprule
\textbf{Rank} & \textbf{Configuration} & \textbf{Mean Rank} & $Q^{AB/F}$ & $N^{AB/F}$ & $MI$ & $SSIM$ & $VIFP$ & $EN$ & $SD$ & $SF$ & $AG$ & \textbf{Lat (ms)} \\
\midrule
\textbf{1} & \textbf{K\_1\_no\_clustering} & \textbf{2.8} & \textbf{0.7107} & \textbf{0.0153} & \textbf{3.819} & \textbf{0.7215} & \textbf{0.4764} & \textbf{6.359} & \textbf{0.174} & \textbf{0.048} & \textbf{0.0138} & \textbf{3063} \\
2 & GFF & 3.4 & 0.7356 & 0.0630 & 3.455 & 0.7016 & 0.5034 & 6.230 & 0.165 & 0.053 & 0.0152 & 2273 \\
3 & K\_2 & 5.2 & 0.6410 & 0.0966 & 3.234 & 0.6627 & 0.3965 & 6.391 & 0.185 & 0.055 & 0.0152 & 12655 \\
4 & Averaging & 6.2 & 0.4128 & 0.0000 & 2.964 & 0.7287 & 0.3570 & 5.784 & 0.109 & 0.029 & 0.0087 & 409 \\
5 & softmax\_tau\_1.0 & 6.6 & 0.4727 & 0.0052 & 2.934 & 0.7309 & 0.3542 & 5.917 & 0.132 & 0.033 & 0.0095 & 2239 \\
6 & gamma\_0.0\_no\_edge\_comp & 6.6 & 0.7187 & 0.0438 & 2.949 & 0.6381 & 0.3679 & 6.407 & 0.186 & 0.057 & 0.0158 & 3413 \\
7 & softmax\_tau\_0.5 & 7.2 & 0.4918 & 0.0140 & 2.933 & 0.7288 & 0.3478 & 5.995 & 0.150 & 0.037 & 0.0101 & 2411 \\
8 & T\_5 & 7.6 & 0.6096 & 0.1302 & 2.986 & 0.6389 & 0.3680 & 6.406 & 0.187 & 0.057 & 0.0157 & 4548 \\
9 & T\_10 & 7.6 & 0.6114 & 0.1296 & 2.967 & 0.6385 & 0.3680 & 6.406 & 0.186 & 0.057 & 0.0158 & 6835 \\
10 & softmax\_tau\_0.1 & 7.8 & 0.6215 & 0.0160 & 2.914 & 0.7128 & 0.3309 & 6.184 & 0.183 & 0.047 & 0.0119 & 1678 \\
11 & guided\_radius\_9 & 9.4 & 0.6022 & 0.1213 & 2.833 & 0.6399 & 0.3553 & 6.392 & 0.182 & 0.055 & 0.0153 & 2175 \\
12 & guided\_radius\_13 & 9.8 & 0.5967 & 0.1155 & 2.766 & 0.6420 & 0.3499 & 6.378 & 0.178 & 0.054 & 0.0150 & 2837 \\
13 & gamma\_0.25 & 11.0 & 0.6459 & 0.1327 & 2.512 & 0.6033 & 0.3324 & 6.463 & 0.191 & 0.076 & 0.0198 & 1613 \\
14 & CADM\_v1\_prototype ($\gamma=0.5$) & 13.6 & 0.4894 & 0.2888 & 2.140 & 0.5444 & 0.2875 & 6.521 & 0.198 & 0.108 & 0.0244 & 4686 \\
\bottomrule
\end{tabular}
\end{center}
\end{table*}

\subsection{Test-Set Benchmark Evaluation}
Table IV presents empirical test-set benchmark scores across modalities.

\begin{table}[htbp]
\caption{Empirical Test-Set Benchmark Evaluation (Mean $\pm$ Std).}
\begin{center}
\begin{tabular}{llcccc}
\toprule
\textbf{Dataset} & \textbf{Method} & $Q^{AB/F}$ & $N^{AB/F}$ & $SSIM$ & \textbf{Lat (ms)} \\
\midrule
\multirow{7}{*}{CT-MRI} 
& Averaging & 0.3949 & 0.0000 & 0.7599 & 0.9 \\
& DWT & 0.5479 & 0.0732 & 0.7364 & 9.4 \\
& GFF & \textbf{0.6453} & 0.0317 & 0.7491 & 27.4 \\
& DenseFuse & 0.5074 & 0.0064 & 0.7665 & 317.8 \\
& U2Fusion & 0.3927 & 0.0213 & 0.7530 & 627.5 \\
& SwinFusion & 0.6167 & 0.0416 & \textbf{0.7786} & 9950.3 \\
& \textbf{CADM\_v2} & \underline{0.6296} & \underline{0.0171} & \underline{0.7675} & 88.2 \\
\midrule
\multirow{5}{*}{PET-MRI}
& Averaging & 0.4453 & 0.0003 & 0.7042 & 3.3 \\
& GFF & \textbf{0.7952} & 0.0149 & 0.7410 & 25.5 \\
& DenseFuse & 0.6475 & \textbf{0.0044} & 0.7427 & 235.5 \\
& SwinFusion & 0.7424 & 0.0393 & \textbf{0.7528} & 11346.8 \\
& \textbf{CADM\_v2} & \underline{0.7277} & \underline{0.0051} & \underline{0.7502} & 92.2 \\
\midrule
\multirow{4}{*}{SPECT-MRI}
& GFF & \textbf{0.7775} & 0.0478 & 0.6701 & 31.0 \\
& DenseFuse & 0.4516 & \textbf{0.0028} & \underline{0.7009} & 277.9 \\
& SwinFusion & 0.6351 & 0.0362 & \textbf{0.7019} & 13177.5 \\
& \textbf{CADM\_v2} & \underline{0.7001} & \underline{0.0257} & 0.6901 & 114.2 \\
\midrule
\multirow{4}{*}{Lytro}
& GFF & \textbf{0.7700} & 0.0688 & 0.8256 & 256.9 \\
& DenseFuse & 0.5220 & \textbf{0.0010} & \textbf{0.8766} & 1115.8 \\
& SwinFusion & \underline{0.7461} & 0.0190 & 0.8565 & 66383.6 \\
& \textbf{CADM\_v2} & 0.7332 & \underline{0.0084} & \underline{0.8465} & 1657.3 \\
\bottomrule
\end{tabular}
\end{center}
\end{table}

\subsection{Statistical Significance}
Table V summarizes paired Wilcoxon signed-rank tests against the closest training-free competitor (GFF).

\begin{table}[htbp]
\caption{Paired Wilcoxon Signed-Rank Tests against GFF.}
\begin{center}
\begin{tabular}{llccc}
\toprule
\textbf{Dataset} & \textbf{Metric} & \textbf{Diff} & \textbf{Wilcoxon p} & \textbf{Holm-Adj p} \\
\midrule
CT-MRI & $Q^{AB/F}$ & -0.0157 & 0.0067 & 0.0134 * \\
CT-MRI & $N^{AB/F}$ & -0.0146 & 0.0001 & 0.0005 *** \\
CT-MRI & $SSIM$ & +0.0184 & 0.0007 & 0.0026 ** \\
CT-MRI & $VIFP$ & -0.0572 & 0.0001 & 0.0004 *** \\
PET-MRI & $Q^{AB/F}$ & -0.0674 & 0.0001 & 0.0004 *** \\
PET-MRI & $N^{AB/F}$ & -0.0098 & 0.0007 & 0.0026 ** \\
PET-MRI & $MI$ & +0.1693 & 0.0001 & 0.0004 *** \\
PET-MRI & $SSIM$ & +0.0093 & 0.0007 & 0.0026 ** \\
SPECT-MRI & $N^{AB/F}$ & -0.0220 & 0.0001 & 0.0004 *** \\
SPECT-MRI & $SSIM$ & +0.0200 & 0.0001 & 0.0004 *** \\
\bottomrule
\end{tabular}
\end{center}
\end{table}

\subsection{Qualitative Analysis}
Test pairs categorized by $Q^{AB/F}$:
(1) CT-MRI: Worst case \texttt{CT-MRI\_16016} ($Q^{AB/F}=0.5978$, $SSIM=0.7889$); Median case \texttt{CT-MRI\_16019} ($Q^{AB/F}=0.6288$, $SSIM=0.8306$); Best case \texttt{CT-MRI\_20005} ($Q^{AB/F}=0.6772$, $SSIM=0.7290$).
(2) PET-MRI: Worst case \texttt{PET-MRI\_25032} ($Q^{AB/F}=0.6858$); Median case \texttt{PET-MRI\_25027} ($Q^{AB/F}=0.7355$); Best case \texttt{PET-MRI\_25020} ($Q^{AB/F}=0.7703$).
(3) SPECT-MRI: Worst case \texttt{SPECT-MRI\_10025} ($Q^{AB/F}=0.6648$); Median case \texttt{SPECT-MRI\_10017} ($Q^{AB/F}=0.6977$); Best case \texttt{SPECT-MRI\_10012} ($Q^{AB/F}=0.7444$).
Dense skull bone in \texttt{CT-MRI\_16016} exhibits slight edge softness relative to GFF due to conservative diffusion conduction ($\kappa=0.1176$), but difference maps verify zero edge-ringing or halo artifacts ($N^{AB/F} = 0.0162$ vs 0.0317 for GFF).

\subsection{Computational Efficiency}
Table VI documents computational runtime and memory profiling on CPU. CADM-MSWN v2 operates in 88.2 ms on CPU, running over 100$\times$ faster than SwinFusion.

\begin{table}[htbp]
\caption{Computational Efficiency Benchmarking (CPU Execution).}
\begin{center}
\begin{tabular}{llccc}
\toprule
\textbf{Method} & \textbf{Type} & \textbf{Params} & \textbf{Latency} & \textbf{Peak RAM} \\
\midrule
Averaging & Classical & 0 & 0.9 ms & 114 MB \\
DWT & Classical & 0 & 9.4 ms & 125 MB \\
GFF \cite{ref_gff} & Classical & 0 & 27.4 ms & 185 MB \\
\textbf{CADM\_v2} & Proposed & \textbf{0} & \textbf{88.2 ms} & \textbf{195 MB} \\
DenseFuse \cite{ref_densefuse} & Deep AE & 73,985 & 317.8 ms & 420 MB \\
U2Fusion \cite{ref_u2fusion} & Deep DenseNet & 659,105 & 627.5 ms & 480 MB \\
SwinFusion \cite{ref_swinfusion} & Vision Trans. & 1,245,840 & 9,950.3 ms & 620 MB \\
\bottomrule
\end{tabular}
\end{center}
\end{table}

\section{Limitations}
\begin{enumerate}
    \item \textbf{Deep Baseline In-Domain Advantage}: U2Fusion, SwinFusion, and CDDFuse were trained directly on RoadScene, MSRS, and Harvard Medical Atlas data, providing in-domain memorization advantages.
    \item \textbf{Metric Sensitivity}: $N^{AB/F}$ naturally favors conservative smoothers. $Q^{AB/F}$ values vary across toolboxes due to angle unwrapping conventions.
    \item \textbf{Hardware Execution}: All benchmarks were executed on CPU for fair comparison; deep models can leverage dedicated GPUs.
    \item \textbf{Spatial Registration}: Perfect spatial co-registration is assumed. Real-world deformation requires prior elastic registration.
    \item \textbf{Statistical Cohort Size}: Medical test cohorts comprise 15 held-out pairs each; statistical power is inherently constrained.
    \item \textbf{Standoff vs GFF}: GFF achieves higher $Q^{AB/F}$ on sharp edges, but injects higher edge noise ($N^{AB/F} \sim 0.03-0.07$). Our method should be viewed as an ablation-driven study establishing what components genuinely matter.
\end{enumerate}

\section{Conclusion}
We presented a training-free multimodal image fusion pipeline and an extensive ablation study showing which components genuinely matter. Our development sweep across 120 pairs showed that edge compensation ($\gamma > 0$) and clustering ($K > 1$) actively harm structural fidelity, whereas a clean additive two-layer decomposition ($GT + FT = Y$) with linear local-statistics weighting achieves the best overall performance. On held-out test data across 7 datasets, the frozen method does not outperform GFF on edge preservation ($Q^{AB/F}$), but produces significantly lower artifact noise ($N^{AB/F}$) and competitive structural similarity with zero training requirements and 88 ms CPU latency. We release the full protocol, configurations, and results.

\section*{Declarations}
\textbf{Funding}: No external funding was received. \textbf{Conflicts of Interest}: None declared. \textbf{Data Availability}: All datasets (MSRS, RoadScene, TNO, Harvard Whole Brain Atlas, Lytro) are publicly available. \textbf{Code Availability}: Code, configurations, split manifests, and raw evaluation tables are publicly released. \textbf{Ethics}: Open-access anonymized medical data were used.

\begin{thebibliography}{00}
\bibitem{ref_pm} P. Perona and J. Malik, ``Scale-space and edge detection using anisotropic diffusion,'' \textit{IEEE Trans. Pattern Anal. Mach. Intell.}, vol. 12, no. 7, pp. 629--639, 1990.
\bibitem{ref_guided} K. He, J. Sun, and X. Tang, ``Guided image filtering,'' \textit{IEEE Trans. Pattern Anal. Mach. Intell.}, vol. 35, no. 6, pp. 1397--1409, 2013.
\bibitem{ref_gff} S. Li, X. Kang, and J. Hu, ``Image fusion with guided filtering,'' \textit{IEEE Trans. Image Process.}, vol. 22, no. 7, pp. 2864--2875, 2013.
\bibitem{ref_bavirisetti} D. P. Bavirisetti and R. Dhuli, ``Fusion of infrared and visible sensor images based on anisotropic diffusion and Karhunen-Loeve transform,'' \textit{IEEE Sensors J.}, vol. 16, no. 1, pp. 203--209, 2016.
\bibitem{ref_cbf} B. K. S. Kumar, ``Image fusion based on pixel significance using cross bilateral filter,'' \textit{Signal Image Video Process.}, vol. 9, no. 5, pp. 1193--1204, 2015.
\bibitem{ref_densefuse} H. Li and X.-J. Wu, ``DenseFuse: A fusion approach to infrared and visible images,'' \textit{IEEE Trans. Image Process.}, vol. 28, no. 5, pp. 2614--2623, 2019.
\bibitem{ref_u2fusion} H. Xu, J. Ma, J. Jiang, X. Guo, and H. Ling, ``U2Fusion: A unified unsupervised image fusion network,'' \textit{IEEE Trans. Pattern Anal. Mach. Intell.}, vol. 44, no. 1, pp. 502--518, 2022.
\bibitem{ref_swinfusion} J. Ma, L. Tang, F. Fan, J. Huang, X. Mei, and Y. Ma, ``SwinFusion: Cross-domain long-range learning for general image fusion via Swin Transformer,'' \textit{IEEE/CAA J. Autom. Sin.}, vol. 9, no. 7, pp. 1205--1217, 2022.
\bibitem{ref_cddfuse} Z. Zhao, H. Bai, J. Zhang, Y. Zhang, S. Xu, Z. Lin, R. Timofte, and L. Van Gool, ``CDDFuse: Correlation-driven dual-branch feature decomposition for multi-modality image fusion,'' in \textit{Proc. CVPR}, 2023.
\bibitem{ref_qabf} C. S. Xydeas and V. Petrovic, ``Objective image fusion performance measure,'' \textit{Electron. Lett.}, vol. 36, no. 4, pp. 308--309, 2000.
\bibitem{ref_nabf} V. Petrovic and C. Xydeas, ``Objective evaluation of hybrid image fusion performance,'' in \textit{Proc. SPIE Defense and Security}, 2005.
\bibitem{ref_vif} H. R. Sheikh and A. C. Bovik, ``Image information and visual quality,'' \textit{IEEE Trans. Image Process.}, vol. 15, no. 2, pp. 430--444, 2006.
\bibitem{ref_viff} Y. Han, Y. Cai, Y. Cao, and X. Xu, ``A new image fusion performance metric based on visual information fidelity,'' \textit{Information Fusion}, vol. 14, no. 2, pp. 127--135, 2013.
\bibitem{ref_ssim} Z. Wang, A. C. Bovik, H. R. Sheikh, and E. P. Simoncelli, ``Image quality assessment: from error visibility to structural similarity,'' \textit{IEEE Trans. Image Process.}, vol. 13, no. 4, pp. 600--612, 2004.
\bibitem{ref_tno} A. Toet, ``The TNO multiband image data collection,'' \textit{Data in Brief}, vol. 15, pp. 249--253, 2017.
\bibitem{ref_piafusion} L. Tang, J. Yuan, H. Zhang, X. Jiang, and J. Ma, ``PIAFusion: A progressive infrared and visible image fusion network based on illumination aware,'' \textit{Information Fusion}, vol. 83-84, pp. 111--122, 2022.
\bibitem{ref_harvard} K. A. Johnson and J. A. Becker, ``The Whole Brain Atlas,'' Harvard Medical School, 1999.
\bibitem{ref_lytro} M. Nejati, S. Samavi, and S. Shirani, ``Multi-focus image fusion using dictionary-based sparse representation,'' \textit{Information Fusion}, vol. 25, pp. 72--84, 2015.
\bibitem{ref_nsct} A. L. da Cunha, J. Zhou, and M. N. Do, ``The nonsubsampled contourlet transform: theory, design, and applications,'' \textit{IEEE Trans. Image Process.}, vol. 15, no. 10, pp. 3089--3101, 2006.
\end{thebibliography}

\end{document}
